% Executive summary (front matter, placed before the table of contents).
\phantomsection
\addcontentsline{toc}{section}{Executive summary}
\section*{Executive summary}

\paragraph{Purpose and status.} This document is an investigation draft for a master's thesis. It states one research idea, analyzes the baseline paper it builds on, and collects the unit designs, learning methods, implementation steps and experiments needed to test the idea. No simulation or hardware result exists yet; every statement about the proposed design is reasoning, and each idea carries the status \emph{active}, \emph{parked} or \emph{rejected} (Section~\ref{sec:ledger}).

\paragraph{The idea.} The hidden layer of a spiking neural network is built from the FPGA's own logic block: a look-up table (LUT), a flip-flop and an output multiplexer, with feedback through the routing fabric (Section~\ref{sec:idea}). This is not an ordinary SNN implemented on an FPGA; the computational unit is made of the resources every logic block already contains. Two variants are considered (Section~\ref{sec:unit}): the minimal unit, a single LUT--FF--MUX block whose four truth-table bits are the learned parameters, and a cluster of several blocks with wider state and higher fan-in. Input and output layers are unrestricted, with delta encoding at the input and counters with argmax at the output; the hidden layer is the research object. Units communicate in one-bit events, hold persistent state, and update in parallel on every clock; any latency advantage over the baseline's serial scan is a design estimate, not a measurement (Section~\ref{sec:impl}).

\paragraph{Expected contribution.} Nothing is claimed yet. A contribution has to be established through a literature review and measured results, plausibly around the FPGA-native recurrent hidden unit, its state-transition structure, its resource and timing behavior, and recurrent trajectory comparison inside that structure (Section~\ref{sec:idea}). The feed-forward part has prior art in LogicNets, NeuraLUT and differentiable logic and weightless networks; the recurrent state-carrying unit with delta or differential-time coding may be the open part.

\paragraph{What the baseline gives and lacks.} The reference paper, Windhager et~al.\ (Section~\ref{sec:paper}), runs a fully connected SNN on MNIST (784--1000--10) at 97.00\% accuracy with 6-bit weights, about 838 images/s and 5.46 mJ/image on a Zynq ZCU102. Its strengths are multiplier-free arithmetic, robustness to coarse weights, and a differential-time code that decouples spike time from processing time. A derived cost model shows the limits: latency grows with fan-in times the number of distinct spike times, so pruning saves nothing; layers are processed one after another over stored spike trains, which rules out recurrence and lateral links; there is one global decay constant; and the evidence rests on a single run on static data, with training--hardware equivalence unverified and resource figures that need checking.

\paragraph{Relation to the LIF neuron.} By the strict definition the logic-block unit is not an LIF neuron (Section~\ref{sec:snnq}); it is a quantized four-state spiking unit, closest to a finite-state machine that emits spikes. What is lost, namely the weighted sum over many synapses, the graded potential with decay and the timing information, can be recovered inside the same primitives (Section~\ref{sec:recover}): a dendrite tree of branch and merge LUTs at about $n/3$ LUT4s per neuron for $n$ inputs, a wider saturating soma with a shared leak tick that replaces the decay constant, and a time replay counter that turns the differential time code into clock time.

\paragraph{Learning.} Plain backpropagation does not apply: truth-table bits have no slope, state is recurrent, and counting followed by argmax is not differentiable (Section~\ref{sec:learning}). The candidate methods are policy gradient over bit probabilities, evolutionary search in a bit-exact simulator, a reward-modulated local rule that would need writable table storage, and enumeration of a trained few-bit neuron, which gives a starting point that search can refine. Differentiable relaxations are set aside for now.

\paragraph{Change detection.} A Siamese variant runs two recurrent paths with shared tables and compares state and transition differences through a third recurrent cluster that keeps its own mismatch history (Section~\ref{sec:siamese}). It targets transient, sustained, repeated and reversed changes between two observations; it is an experimental possibility, not a proven capability.

\paragraph{Path to results.} The build order is: reproduce the baseline in snnTorch, write a bit-exact software model of the logic-block network, implement the unit in SystemVerilog and verify it against the model, add the layer structure, then load to the board and measure (Section~\ref{sec:impl}). The board is a Sipeed Tang Primer 20K; the first target is a sparse network, with the image downsampled to $14\times14$, about 100 hidden units with fan-in 16--64, and 10 output counters. The experimental plan (Section~\ref{sec:experiments}) covers temporal-pattern discrimination, synthetic change detection with paired paths, an MNIST comparison at an equal resource budget, and scaling and ablation studies; an improvement is attributed to a feature only if an ablation isolates it.

\paragraph{Main risks.} The open questions (Section~\ref{sec:open}) reduce to five risks. (i) The advantage over finite-state machines, reservoir computing and weightless networks is unproven, and the novelty claim is not yet established. (ii) A two-bit state remembers only the last change, so long dependencies need depth or lateral loops. (iii) Fan-in four and the dendrite-tree cost limit connectivity, and the right sparsity level is unknown. (iv) No candidate training method is known to scale to the number of table bits. (v) Recurrence forces clock-locked updates, which may dissolve the baseline's decoupling of spike time and processing time.

\paragraph{Idea ledger snapshot.} Of the 21 ideas collected so far, 12 are active, 7 are parked and 2 are rejected: quantum computing, and reliance on physical effects such as metastability or RC elements. Each entry carries pros, cons and a next test (Section~\ref{sec:ledger}).

\newpage